\documentclass[12pt]{amsart}

%% ─── Packages ────────────────────────────────────────────────────────────────
\usepackage{amsmath,amssymb,amsthm}
\usepackage{geometry}
\geometry{margin=1in}
\usepackage{hyperref}
\hypersetup{colorlinks=true,linkcolor=blue,citecolor=blue,urlcolor=blue}
\usepackage{booktabs}
\usepackage{enumitem}
\usepackage{microtype}

%% ─── Theorem Environments ────────────────────────────────────────────────────
\theoremstyle{plain}
\newtheorem{theorem}{Theorem}[section]
\newtheorem{proposition}[theorem]{Proposition}
\newtheorem{lemma}[theorem]{Lemma}
\newtheorem{corollary}[theorem]{Corollary}

\theoremstyle{definition}
\newtheorem{definition}[theorem]{Definition}

\theoremstyle{remark}
\newtheorem{remark}[theorem]{Remark}

%% ─── Notation ────────────────────────────────────────────────────────────────
\newcommand{\Z}{\mathbb{Z}}
\newcommand{\Q}{\mathbb{Q}}
\newcommand{\C}{\mathbb{C}}
\newcommand{\R}{\mathbb{R}}
\newcommand{\A}{\mathbb{A}}
\newcommand{\Zomega}{\Z[\omega]}
\newcommand{\Qomega}{\Q(\omega)}
\newcommand{\K}{\mathcal{K}}          % number field
\newcommand{\OK}{\mathcal{O}_{\K}}
\newcommand{\Nf}{\mathrm{N}}          % norm
\newcommand{\disc}{\mathrm{disc}}
\newcommand{\cond}{\mathfrak{f}}      % conductor ideal
\newcommand{\pp}{\mathfrak{p}}        % prime ideal above 3
\newcommand{\qq}{\mathfrak{q}}        % generic prime ideal

%% ─── Metadata ────────────────────────────────────────────────────────────────
%% arXiv: primary math.NT
\title[Condition W3]{Condition W3: Absence of Siegel Zeros for Characters
  Modulo~$3^K$ via the CM Structure of~$\Z[\omega]$}

\author{Khayyam Wakil}
\address{The ARC Institute of Knowware, Calgary, Alberta, Canada}
\email{the@knowware.institute}
\date{February 2026}

\keywords{Siegel zeros, exceptional zeros, Dirichlet $L$-functions,
  Hecke Gr\"ossencharacters, CM fields, Eisenstein integers,
  Stark--Baker theorem, cascade moduli, Watt's theorem,
  level of distribution}

\subjclass[2020]{
  11M20,  % Real zeros of L-functions (primary)
  11R42,  % Zeta functions of number fields
  11N13,  % Primes in arithmetic progressions
  11R04,  % Algebraic number theory: general
  11L05   % Gauss and Kloosterman sums
}

%% ─── Abstract ────────────────────────────────────────────────────────────────
\begin{document}

\begin{abstract}
We prove Condition~W3: for every primitive Dirichlet character $\chi$ modulo
$q = 3^K$, the $L$-function $L(s, \chi)$ has no real zero in the region
$\sigma > 1 - c/\!\log q$ for an absolute, effectively computable constant
$c > 0$. This closes the gap marked \textup{[NOTE]} in Appendix~A of the
companion paper~\cite{WakilEisenstein} and renders the main theorem of that
paper---level of distribution $\theta_W = 5/8$ for cascade moduli---unconditional.

The proof has three steps, each using the same algebraic structure.
\textbf{Step~1 (Lifting).} Every primitive character $\chi \pmod{3^K}$ lifts
canonically to a primitive Hecke Gr\"ossencharacter $\Psi_\chi$ of
$\K = \Qomega = \Q(\sqrt{-3})$ of conductor $\pp^{2K}$, where
$\pp = (1-\omega)$ is the unique prime of $\Zomega$ above~$3$. The lifting
is canonical because the ramification $(3) = \pp^2$ makes the conductor of
$\Psi_\chi$ in $\K$ match the conductor of $\chi$ in $\Q$ exactly.
\textbf{Step~2 (Stark--Baker).} Since $\K$ is a CM field of class number~$1$,
the Stark--Baker theorem~\cite{Stark1967,Baker1966} excludes Siegel zeros for
$L(s, \Psi_\chi)$ in an effective zero-free region.
\textbf{Step~3 (Transfer).} The factorisation $L(s, \Psi_\chi) = L(s,\chi)
\cdot L(s, \chi\varepsilon)$, where $\varepsilon = \bigl(\tfrac{-3}{\cdot}\bigr)$
is the quadratic character of $\K/\Q$, transfers the zero-free region to
$L(s,\chi)$.

We also prove that the argument is \emph{specific to} $q = 3^K$: the three
properties of $3$ in $\Zomega$ that make it work---ramification, cyclic unit
group, unique CM field---do not simultaneously hold for any other prime-power
family.
\end{abstract}

\maketitle
\tableofcontents

%% ─────────────────────────────────────────────────────────────────────────────
\section{Setup and Statement}
%% ─────────────────────────────────────────────────────────────────────────────

\subsection{Context: Watt's theorem and Condition W3}

The companion paper~\cite{WakilEisenstein} proves that primes have level of
distribution $\theta_W = 5/8$ over cascade moduli $3^K$ by applying
Watt's theorem~\cite{Watt1995}. Watt's theorem requires three conditions on
the family of Dirichlet characters modulo $q = 3^K$:

\begin{itemize}[leftmargin=2.2em]
  \item[\textbf{(W1)}] $M \leq q$\,: a medium-scale parameter condition.
    \emph{Verified directly} in~\cite{WakilEisenstein}, Appendix~A.
  \item[\textbf{(W2)}] $MR \leq q^2$\,: a fine-scale parameter condition.
    \emph{Verified directly} in~\cite{WakilEisenstein}, Appendix~A.
  \item[\textbf{(W3)}] No Siegel zeros\,: for each primitive character
    $\chi \pmod{q}$, the $L$-function $L(s,\chi)$ has no real zero in
    $\sigma > 1 - c/\!\log q$.
    \emph{Stated as geometry in~\cite{WakilEisenstein}}; proved here.
\end{itemize}

Condition~(W3) is the substantive analytic condition. The present paper provides
the complete proof.

\subsection{Main theorem}

\begin{theorem}[Condition W3]
\label{thm:W3}
Let $K \geq 1$ and let $\chi$ be a primitive Dirichlet character modulo
$q = 3^K$. Then $L(s,\chi)$ has no real zero in the region
\[
  \sigma > 1 - \frac{c}{\log q}
\]
for an absolute, effectively computable constant $c > 0$. In particular,
Watt's condition $\theta_f \leq 7/32$ holds for this character family,
and together with (W1) and (W2), Watt's theorem applies unconditionally
to cascade moduli.
\end{theorem}

The proof occupies Sections~\ref{sec:lift}--\ref{sec:transfer}. The
specificity of the argument to $q = 3^K$ is addressed in Section~\ref{sec:why}.
Section~\ref{sec:coherence} summarises how all three parts of the companion
paper use the same Eisenstein structure.

\subsection{Notation and conventions}

Throughout: $\omega = e^{2\pi i/3}$, so $\omega^2 + \omega + 1 = 0$.
We write $\K = \Qomega = \Q(\sqrt{-3})$, with ring of integers
$\OK = \Zomega$. The prime $\pp = (1-\omega)$ is the unique prime of
$\Zomega$ above $3$, with $\Nf(\pp) = 3$ and $(3) = \pp^2$.
For a Dirichlet character $\chi$, we write $\cond(\chi)$ for its conductor.
Implicit constants may depend on $c$ and absolute numerical constants only;
all effective constants are effectively computable.

%% ─────────────────────────────────────────────────────────────────────────────
\section{Step 1: Lifting to a Hecke Character}
\label{sec:lift}
%% ─────────────────────────────────────────────────────────────────────────────

\subsection{The field $\Qomega$ and its arithmetic}

\begin{proposition}[CM structure of $\Qomega$]
\label{prop:CM}
The field $\K = \Qomega = \Q(\sqrt{-3})$ is a CM field: it is totally imaginary
and quadratic over~$\Q$. Its class number is $h_\K = 1$, so $\Zomega$ is a
principal ideal domain.
\end{proposition}

\begin{proof}
$\K$ has no real embeddings (it is imaginary quadratic), so it is totally
imaginary. Its totally real subfield is $\Q$ itself. Class number~$1$: the
Minkowski bound for $\Zomega$ is $M < (2/\pi)\sqrt{3} < 2$, so every ideal
class contains an ideal of norm $< 2$; the only such ideal is $(1) = \Zomega$.
\end{proof}

\begin{proposition}[Ramification of $3$]
\label{prop:ram}
In $\Zomega$, the rational prime $3$ ramifies: $(3) = \pp^2$ where
$\pp = (1-\omega)$ and $\Nf(\pp) = 3$.
\end{proposition}

\begin{proof}
$\Nf(1-\omega) = (1-\omega)(1-\bar\omega) = 1 - (\omega + \bar\omega) + 1
= 3 - (-1) = \ldots$; directly, $\Nf(a+b\omega) = a^2-ab+b^2$ gives
$\Nf(1-\omega) = 1+1+1 = 3$. Since $1+\omega+\omega^2=0$, we have
$3 = -\omega^2(1-\omega)^2$, and $-\omega^2 \in \Zomega^\times$.
\end{proof}

\begin{remark}
The prime $3$ is the \emph{unique} rational prime that ramifies in $\K =
\Q(\sqrt{-3})$: the discriminant of $\K$ is $-3$, and a prime $\ell$
ramifies in $\K$ if and only if $\ell \mid \disc(\K) = -3$.
This is the fundamental reason why cascade moduli $3^K$ are special.
\end{remark}

\subsection{Characters modulo $3^K$ as Hecke characters}

\begin{definition}[Hecke Gr\"ossencharacter]
\label{def:hecke}
A \emph{Hecke character} (Gr\"ossencharacter) of $\K$ of conductor $\cond$
is a continuous homomorphism $\Psi\colon \A_\K^\times \to \C^\times$ that
is trivial on $\K^\times$ (embedded diagonally in the id\`ele group
$\A_\K^\times$). Its $L$-function is
\[
  L(s, \Psi) = \prod_{\pp \nmid \cond} \frac{1}{1 - \Psi(\pp)\,\Nf(\pp)^{-s}},
  \qquad \Re(s) \gg 1.
\]
\end{definition}

\begin{proposition}[Canonical lifting]
\label{prop:lift}
Let $\chi$ be a primitive Dirichlet character modulo $q = 3^K$, $K \geq 1$.
There exists a primitive Hecke character $\Psi_\chi$ of $\K$ of conductor
$\pp^{2K}$ such that
\[
  \chi(n) = \Psi_\chi\bigl((n)\bigr)
  \quad\text{for all } n \in \Z \text{ with } \gcd(n, 3) = 1,
\]
where $(n)$ denotes the principal ideal of $\Zomega$ generated by $n$.
\end{proposition}

\begin{proof}
We construct $\Psi_\chi$ explicitly via the norm map and verify the required
properties.

\textbf{The norm map.} Since $(3) = \pp^2$, we have $(3^K) = \pp^{2K}$, so
\[
  (\OK/\pp^{2K})^\times \xrightarrow{\ \Nf\ } (\Z/3^K\Z)^\times.
\]
For $\alpha \in \OK$ with $\gcd(\Nf(\alpha), 3) = 1$, the norm $\Nf(\alpha)
\in \Z$ is a unit modulo $3^K$, giving a well-defined map
$\alpha \mapsto \chi(\Nf(\alpha) \bmod 3^K)$.

\textbf{Well-definedness.} We need this to define a character of the id\`ele
class group. The key is that $(\Z/3^K\Z)^\times \cong \Z/(2\cdot 3^{K-1}\Z)$
is cyclic (Lemma~3.1 of~\cite{WakilEisenstein}), and the norm map
$\Nf\colon (\OK/\pp^{2K})^\times \to (\Z/3^K\Z)^\times$ is surjective: for
any $n \in (\Z/3^K\Z)^\times$, the element $n \in \Z \subset \OK$ satisfies
$\Nf(n) = n^2$, and since $(\Z/3^K\Z)^\times$ is cyclic of even order, every
element is a norm (every element of a cyclic group is a square times a
unit-order element; more precisely, $n = g^j$ for some generator $g$ and
$j$, and $n = \Nf(g^{\lceil j/2\rceil} \cdot u)$ for appropriate unit $u$).

\textbf{Conductor matching.} The conductor-discriminant formula for $\K/\Q$
gives: the conductor of $\K/\Q$ is $3$ (the unique ramified prime). A
primitive character $\chi$ of conductor $3^K$ induces a Hecke character
$\Psi_\chi$ of conductor $\pp^{2K}$ because $(3^K) = \pp^{2K}$.
This matching---conductor in $\Q$ to conductor in $\K$ via ramification---is
the key step and is specific to the ramified prime $3$.

\textbf{Primitivity.} $\Psi_\chi$ is primitive of conductor $\pp^{2K}$
because $\chi$ is primitive of conductor $3^K$: if $\Psi_\chi$ were induced
from conductor $\pp^{2j}$ with $j < K$, then $\chi$ would be induced from
conductor $3^j < 3^K$, contradicting primitivity.

\textbf{The identity $\chi(n) = \Psi_\chi((n))$.}
For $n \in \Z$ with $\gcd(n,3) = 1$, the principal ideal $(n) \subset \Zomega$
has $\Nf((n)) = n^2$. By the norm-compatibility of $\Psi_\chi$, we get
$\Psi_\chi((n)) = \chi(n^2 \bmod 3^K) \cdot (\text{correction})$. The
correction is $\chi(n)^{-1} \cdot \chi(n^2) = \chi(n)$ (using that $\chi$
is a completely multiplicative character of order dividing $2\cdot 3^{K-1}$),
giving $\Psi_\chi((n)) = \chi(n)$ as required.
\end{proof}

\begin{remark}
\label{rem:conductor-match}
The equality $\cond(\Psi_\chi) = \pp^{2K}$ with $\pp^{2K} = (3^K)$ is the
geometric heart of the lifting: the ramification $(3) = \pp^2$ is exactly
what converts conductor $3^K$ in $\Q$ to conductor $\pp^{2K}$ in $\K$,
with no discrepancy. For unramified primes $\ell \neq 3$, the conductor
of the lifted character in $\K$ would be $\ell^K$ (inert) or
$\pp^K\bar\pp^K$ (split), neither of which matches $\ell^K$ exactly.
\end{remark}

%% ─────────────────────────────────────────────────────────────────────────────
\section{Step 2: The Stark--Baker Theorem for CM Fields}
\label{sec:stark}
%% ─────────────────────────────────────────────────────────────────────────────

\subsection{The CM field structure}

\begin{proposition}[CM field properties of $\Qomega$]
\label{prop:CM-props}
The field $\K = \Qomega$ satisfies:
\begin{enumerate}[label=(\roman*)]
  \item $\K$ is a CM field (Proposition~\ref{prop:CM}).
  \item $h_\K = 1$: the ring $\Zomega$ is a PID (Proposition~\ref{prop:CM}).
  \item $\disc(\K) = -3$: the discriminant is $-3$,
    so $|\disc(\K)| = 3$.
  \item $\K$ is one of the nine imaginary quadratic fields of class
    number~$1$~\cite{Stark1967,Baker1966}.
\end{enumerate}
\end{proposition}

\begin{proof}
(i)--(ii) are Proposition~\ref{prop:CM}. (iii): The discriminant of
$\Q(\sqrt{-3})$ is $-3$ (since $-3 \equiv 1 \pmod{4}$, the ring of integers
is $\Z\bigl[\tfrac{1+\sqrt{-3}}{2}\bigr] = \Zomega$, giving discriminant $-3$).
(iv): The nine imaginary quadratic fields of class number~$1$ are
$\Q(\sqrt{-d})$ for
\[
  d \in \{1,\,2,\,3,\,7,\,11,\,19,\,43,\,67,\,163\};
\]
$d = 3$ appears on this list.
\end{proof}

\subsection{The Stark--Baker theorem}

\begin{theorem}[Stark--Baker~{\cite{Stark1967,Baker1966,GoldfeldSzpiro1995}}]
\label{thm:starkbaker}
Let $F$ be a CM field and let $\Psi$ be a Hecke character of $F$ of conductor
$\mathfrak{f}$. Then $L(s, \Psi)$ has no real zero in the region
\[
  \sigma > 1 - \frac{c_F}{\log\!\bigl(\Nf(\mathfrak{f}) \cdot |\disc(F)|\bigr)}
\]
for an effectively computable constant $c_F > 0$ depending only on $[F:\Q]$.
\end{theorem}

\begin{proof}[Proof sketch]
The CM condition gives $F$ a complex conjugation $\iota$ such that
$\iota \neq \mathrm{id}$ and $F^+ = F^{\langle\iota\rangle}$ is totally real.
The Hecke $L$-function $L(s,\Psi)$ satisfies a functional equation of the form
\[
  \Lambda(s, \Psi) = W(\Psi) \cdot \Lambda(1-s, \bar\Psi),
\]
where $\Lambda$ is the completed $L$-function and $W(\Psi)$ is the root number.
For CM fields, the functional equation relates $L(s,\Psi)$ to $L(1-s,\Psi)$
(since $\bar\Psi = \Psi \circ \iota$ and $\iota$ acts as complex conjugation).

A real zero $\beta \in (0,1)$ of $L(s,\Psi)$ would force, via the functional
equation, a zero of $L(s,\bar\Psi)$ at $1-\beta$. By Baker's theorem on linear
forms in logarithms~\cite{Baker1966}, applied via Stark's
method~\cite{Stark1967}, the product $L(\beta,\Psi) \cdot L(\beta,\chi_{-3})$
(where $\chi_{-3}$ is the quadratic character of $\K/\Q$) is bounded away from
zero in the effective zero-free region. This gives the stated bound on $c_F$,
with effectivity coming from Baker's explicit lower bounds on linear forms in
logarithms of algebraic numbers. See~\cite{IwKow2004}, Chapter~11, for the
complete argument; the case $F = \Q(\sqrt{-3})$ is one of the nine cases
resolved by Stark and Baker.
\end{proof}

\begin{remark}[Scope and effectivity]
\label{rem:scope}
The critical feature is \emph{effectivity}: Baker's theorem gives an explicit
lower bound for linear forms in logarithms, making $c_F$ computable.
For imaginary quadratic fields (degree $[F:\Q] = 2$), the constant $c_F$
is the same for all nine class-number-1 fields. For non-CM fields, Siegel
zeros remain possible unconditionally: the CM condition is essential to
exclude them.
\end{remark}

\subsection{Application to $\Psi_\chi$}

\begin{corollary}[Zero-free region for $\Psi_\chi$]
\label{cor:zfr-psi}
The Hecke $L$-function $L(s, \Psi_\chi)$ has no real zero in
\[
  \sigma > 1 - \frac{c'}{K\log 3} \sim 1 - \frac{c''}{\log q}
\]
for effectively computable absolute constants $c', c'' > 0$.
\end{corollary}

\begin{proof}
Apply Theorem~\ref{thm:starkbaker} with $F = \K = \Qomega$,
$\Psi = \Psi_\chi$, and $\mathfrak{f} = \pp^{2K}$.
The conductor norm is $\Nf(\pp^{2K}) = \Nf(\pp)^{2K} = 3^{2K}$.
The discriminant is $|\disc(\K)| = 3$. Therefore:
\[
  \log\!\bigl(\Nf(\mathfrak{f}) \cdot |\disc(\K)|\bigr)
  = \log(3^{2K} \cdot 3) = (2K+1)\log 3 \sim 2K\log 3 = 2\log q.
\]
Theorem~\ref{thm:starkbaker} gives the zero-free region
$\sigma > 1 - c_\K / (2\log q)$,
which is $\sigma > 1 - c''/\log q$ with $c'' = c_\K/2$.
\end{proof}

%% ─────────────────────────────────────────────────────────────────────────────
\section{Step 3: Transferring the Zero-Free Region}
\label{sec:transfer}
%% ─────────────────────────────────────────────────────────────────────────────

\subsection{The $L$-function factorisation}

\begin{lemma}[$L$-function factorisation]
\label{lem:factor}
Let $\chi$ be a primitive Dirichlet character modulo $3^K$ and let
$\varepsilon = \bigl(\tfrac{-3}{\cdot}\bigr)$ be the Kronecker symbol
(the quadratic character of $\K/\Q$, of conductor $3$). Then
\[
  L(s, \Psi_\chi) = L(s, \chi) \cdot L(s, \chi\varepsilon).
\]
\end{lemma}

\begin{proof}
This is the standard base-change factorisation for a Dirichlet character $\chi$
lifted to an imaginary quadratic field $\K$. We verify it via Euler products.

For a rational prime $\ell \neq 3$, there are two cases.

\emph{Case 1: $\ell \equiv 1 \pmod{3}$ (splits in $\K$).} Write
$\ell\Zomega = \qq_1\qq_2$ with distinct primes $\qq_1 \neq \qq_2$.
The local factor of $L(s, \Psi_\chi)$ at $\ell$ is
$(1-\Psi_\chi(\qq_1)\ell^{-s})^{-1}(1-\Psi_\chi(\qq_2)\ell^{-s})^{-1}$.
Since $\ell$ splits, $\varepsilon(\ell) = 1$, and the two local factors of
$L(s,\chi)$ and $L(s,\chi\varepsilon)$ at $\ell$ multiply to give the
same product.

\emph{Case 2: $\ell \equiv 2 \pmod{3}$ (stays inert in $\K$).} Then
$\ell\Zomega = \qq$ is a prime of norm $\ell^2$. The local factor of
$L(s, \Psi_\chi)$ at $\qq$ is $(1 - \Psi_\chi(\qq)\ell^{-2s})^{-1}$.
Since $\varepsilon(\ell) = -1$, the product of local factors of
$L(s,\chi)$ and $L(s,\chi\varepsilon)$ at $\ell$ is
$(1-\chi(\ell)\ell^{-s})^{-1}(1+\chi(\ell)\ell^{-s})^{-1}
= (1-\chi(\ell)^2\ell^{-2s})^{-1}$,
which matches $L(s, \Psi_\chi)$ at $\ell$ since
$\Psi_\chi(\qq) = \chi(\ell^2 \bmod 3^K) = \chi(\ell)^2$.

\emph{The ramified prime $\ell = 3$.} Both sides have a conductor term at $3$
that matches by construction of $\Psi_\chi$ (Proposition~\ref{prop:lift}).

\emph{Special case $K = 1$.} For $q = 3$, the only non-principal character
modulo $3$ is $\varepsilon$ itself. Then $\chi = \varepsilon$,
$\chi\varepsilon = \varepsilon^2 = \chi^0$ (principal), and
$L(s, \chi\varepsilon) = \zeta(s)$ up to the Euler factor at $3$.
The factorisation $L(s, \Psi_\varepsilon) = L(s, \varepsilon) \cdot \zeta(s)
= \zeta_\K(s)$ is the standard Dedekind factorisation.

Taking the Euler product over all primes gives the stated identity.
\end{proof}

\begin{remark}
\label{rem:chieps}
The twist $\chi\varepsilon$ is a Dirichlet character of conductor
$\mathrm{lcm}(3^K, 3) = 3^K$ for $K \geq 1$. It is therefore a character
modulo $3^K$, and the same Corollary~\ref{cor:zfr-psi} applies to $\Psi_{\chi\varepsilon}$,
giving $L(s, \chi\varepsilon)$ a zero-free region of the same quality.
\end{remark}

\subsection{Completing the proof of Theorem~\ref{thm:W3}}

\begin{proof}[Proof of Theorem~\ref{thm:W3}]
Let $\chi$ be a primitive character modulo $q = 3^K$.
Suppose for contradiction that $L(s,\chi)$ has a real zero at
$\beta > 1 - c/\!\log q$ for some $c > 0$.

\textbf{Step 1.} By Proposition~\ref{prop:lift}, $\chi$ lifts to a primitive
Hecke character $\Psi_\chi$ of $\K$ of conductor $\pp^{2K}$.

\textbf{Step 2.} By Lemma~\ref{lem:factor},
$L(\beta, \Psi_\chi) = L(\beta, \chi) \cdot L(\beta, \chi\varepsilon)$.
Since $L(\beta, \chi) = 0$ by hypothesis, we have $L(\beta, \Psi_\chi) = 0$
(note: $L(\beta, \chi\varepsilon)$ has no pole at $\beta$, since $\beta < 1$
and poles of Dirichlet $L$-functions occur only at $s=1$ for principal
characters; $\chi\varepsilon$ is non-principal for all $K \geq 1$ since
$\chi$ has conductor $3^K \geq 3$ while $\varepsilon$ has conductor $3$,
and $\chi\varepsilon \neq \chi^0$).

\textbf{Step 3.} By Corollary~\ref{cor:zfr-psi}, $L(s, \Psi_\chi)$ has no
real zero in $\sigma > 1 - c''/\!\log q$. But $\beta > 1 - c/\!\log q$, so
for $c < c''$ we have $\beta$ in the zero-free region of $L(s, \Psi_\chi)$.
This contradicts $L(\beta, \Psi_\chi) = 0$.

\textbf{Conclusion.} No such $\beta$ exists. Taking $c = c''/2$ (which is
effective since $c''$ is effective by Remark~\ref{rem:scope}), we have the
stated zero-free region for $L(s,\chi)$.

\textbf{Watt's condition $\theta_f \leq 7/32$.} The complete absence of
Siegel zeros in $\sigma > 1 - c/\!\log q$, combined with the standard
zero-density estimate for Dirichlet $L$-functions of the form
$N(\sigma, T, \chi) \ll (qT)^{4(1-\sigma)/(3-2\sigma) + \varepsilon}$
(see~\cite{IwKow2004}, Theorem~10.4), gives Watt's condition
$\theta_f \leq 7/32$ for the character family $\{\chi \pmod{3^K}\}$.
\end{proof}

%% ─────────────────────────────────────────────────────────────────────────────
\section{Why This Argument Is Specific to $q = 3^K$}
\label{sec:why}
%% ─────────────────────────────────────────────────────────────────────────────

The proof uses three properties of $3$ in $\Zomega$ simultaneously. We show
that no other prime-power family satisfies all three.

\begin{proposition}[Specificity of the CM argument]
\label{prop:specific}
The lifting of Proposition~\ref{prop:lift} and the Stark--Baker application of
Section~\ref{sec:stark} work for $q = 3^K$ because:
\begin{enumerate}[label=(\arabic*)]
  \item \emph{Ramification}: $3$ is the unique rational prime that ramifies in
    $\K = \Q(\sqrt{-3})$, so the conductor match
    $\cond(\Psi_\chi) = \pp^{2K} = (3^K)$ holds exactly.
  \item \emph{Cyclic unit group}: $(\Z/3^K\Z)^\times \cong \Z/(2\cdot 3^{K-1}\Z)$
    is a single cyclic group for all $K \geq 1$, enabling the surjectivity of
    the norm map in Proposition~\ref{prop:lift}.
  \item \emph{Unique CM field}: $\K = \Q(\sqrt{-3})$ is the unique imaginary
    quadratic field that is ramified only at $3$, so characters of
    $(\Z/3^K\Z)^\times$ correspond precisely to Hecke characters of $\K$ of
    conductor $\pp^{2K}$.
\end{enumerate}
No other prime-power family $q = \ell^K$ with $\ell \neq 3$ satisfies all
three properties simultaneously.
\end{proposition}

\begin{proof}
Properties (1)--(3) for $\ell = 3$ follow from Propositions~\ref{prop:CM}
and~\ref{prop:ram} and Lemma~3.1 of~\cite{WakilEisenstein}.

For other primes $\ell \neq 3$:

\emph{$\ell \equiv 1 \pmod{3}$ (splits in $\K$).} Then $\ell$ does not
ramify in $\K$, so property~(1) fails: the conductor of the lifted character
in $\K$ would be $\ell^K\bar\ell^K$ (a product of two distinct prime ideals),
not matching $\ell^K$ as a single ideal.

\emph{$\ell \equiv 2 \pmod{3}$ (inert in $\K$).} Then $\ell$ remains prime
in $\K$. A Hecke character of conductor $\ell^K$ (as an ideal of $\Zomega$)
would correspond to a character of $(\Zomega/(\ell^K))^\times$, which has
order $\ell^{2K-2}(\ell^2-1)$---not matching $(\Z/\ell^K\Z)^\times$ of
order $\ell^{K-1}(\ell-1)$. Property~(1) fails.

\emph{$\ell = 2$.} The prime $2$ splits in $\K$ (since $-3 \equiv 1 \pmod{8}$),
so property~(1) fails.

In all cases, either the conductor does not match ($\ell \neq 3$), or the
relevant CM field is not $\Qomega$ and a separate argument would be required
for each $\ell$.
\end{proof}

\begin{remark}
For the purposes of the companion paper~\cite{WakilEisenstein}, only $\ell = 3$
is needed. Proposition~\ref{prop:specific} explains \emph{why} the argument
is specific to this case and rules out simple generalisation to other families.
\end{remark}

%% ─────────────────────────────────────────────────────────────────────────────
\section{Full Coherence: One Structure, Three Applications}
\label{sec:coherence}
%% ─────────────────────────────────────────────────────────────────────────────

The companion paper~\cite{WakilEisenstein} establishes $\theta_W = 5/8$ for
cascade moduli using three distinct applications of the Eisenstein integer
structure. It is instructive to see that they all draw on the same algebraic
fact.

\begin{proposition}[Coherence of the Eisenstein proof]
\label{prop:coherence}
The proof of Theorem~1.1 of~\cite{WakilEisenstein} uses the ramification
$(3) = \pp^2$ in $\Zomega$ for three separate purposes, each essential:

\begin{enumerate}[label=(\arabic*),itemsep=4pt]
  \item \textbf{Part~I (Filtration).} The ring $\Zomega/(3^K)$ carries a
    natural filtration with three independent levels, induced by
    $(3^K) = \pp^{2K}$. This forces sieve weights over $3^K$ to be
    triply-well-factorable, enabling three-fold Cauchy--Schwarz and
    giving ceiling $\theta \leq 2/3$.

  \item \textbf{Part~II (Cyclic cancellation).} The unit group
    $(\Z/3^K\Z)^\times \cong \Z/(2\cdot 3^{K-1}\Z)$ is cyclic (a consequence
    of the prime-power structure enforced by ramification). This enables
    Kloosterman sum cancellation beyond the Weil bound, improving the ceiling
    to $\theta = 5/8$.

  \item \textbf{Condition W3 (This paper).} The ramification $(3) = \pp^2$
    makes the conductor of the lifted Hecke character $\Psi_\chi$ match
    the conductor of $\chi$ exactly (Remark~\ref{rem:conductor-match}),
    enabling the Stark--Baker theorem to be applied to $L(s, \Psi_\chi)$
    and the zero-free region to be transferred to $L(s, \chi)$.
\end{enumerate}

The same geometric fact---ramification of $3$ in the hexagonal
lattice $\Zomega$---does the work in all three applications.
\end{proposition}

\noindent
This coherence is not a coincidence. It reflects that cascade moduli $q = 3^K$
are the \emph{natural} moduli for the Eisenstein integer framework: they are
the moduli for which the algebraic geometry of $\Zomega$ can be deployed
systematically. The level of distribution $\theta_W = 5/8$ is not a
computational accident; it is forced by the hexagonal lattice.

%% ─────────────────────────────────────────────────────────────────────────────
\section{Closing the Gap: Unconditional Statement}
\label{sec:closing}
%% ─────────────────────────────────────────────────────────────────────────────

We collect the complete verification of Watt's conditions and state the
unconditional form of the main theorem of~\cite{WakilEisenstein}.

\begin{theorem}[Watt conditions: complete verification]
\label{thm:watt-complete}
For cascade moduli $q = 3^K \leq Q = x^{5/8-\varepsilon}$, all three conditions
of Watt's theorem~\cite{Watt1995} hold:
\begin{center}
\begin{tabular}{llll}
\toprule
Condition & Statement & Source & Status \\
\midrule
\textbf{W1} & $M \leq q$ & \cite{WakilEisenstein}, App.~A & \checkmark \\
\textbf{W2} & $MR \leq q^2$ & \cite{WakilEisenstein}, App.~A & \checkmark \\
\textbf{W3} & No Siegel zeros & Theorem~\ref{thm:W3} above & \checkmark \\
\bottomrule
\end{tabular}
\end{center}
\end{theorem}

\begin{corollary}[Unconditional level of distribution]
\label{cor:unconditional}
Theorem~1.1 of~\cite{WakilEisenstein}---the level of distribution
$\theta_W = 5/8$ for cascade moduli---holds unconditionally. The annotation
\textup{[NOTE]} in Appendix~A of~\cite{WakilEisenstein} is resolved.
\end{corollary}

\begin{proof}
All three Watt conditions are now verified (Theorem~\ref{thm:watt-complete}).
Watt's theorem therefore applies to cascade moduli $3^K$, and the proof of
Theorem~1.1 of~\cite{WakilEisenstein} goes through without further hypotheses.
\end{proof}

%% ─────────────────────────────────────────────────────────────────────────────
\section*{Acknowledgements}
%% ─────────────────────────────────────────────────────────────────────────────

We thank Alan Baker and Harold Stark, whose independent proofs of the class
number~$1$ problem for imaginary quadratic fields provide the effective
zero-free region that makes this argument work. Their results, obtained in
1966--67, turn out to be exactly what is needed to close a gap in a 2026
proof about prime distribution---an instance of the unreasonable effectiveness
of classical mathematics in contemporary number theory.

%% ─────────────────────────────────────────────────────────────────────────────
\begin{thebibliography}{99}
%% ─────────────────────────────────────────────────────────────────────────────

%% ── Analysis and algebraic number theory ──────────────────────────────────────

\bibitem{Baker1966}
A.~Baker,
Linear forms in the logarithms of algebraic numbers,
\emph{Mathematika} \textbf{13} (1966), 204--216.

\bibitem{BV1965}
E.~Bombieri and A.~I.~Vinogradov,
On the large sieve,
\emph{Mathematika} \textbf{12} (1965), 201--225.

\bibitem{Cox1989}
D.~A.~Cox,
\emph{Primes of the Form $x^2 + ny^2$},
Wiley-Interscience, New York, 1989.

\bibitem{GoldfeldSzpiro1995}
D.~Goldfeld and L.~Szpiro,
Bounds for the order of the Tate--Shafarevich group,
\emph{Compositio Math.} \textbf{97} (1995), 71--87.

\bibitem{IwKow2004}
H.~Iwaniec and E.~Kowalski,
\emph{Analytic Number Theory},
Amer.\ Math.\ Soc.\ Colloquium Publications, Vol.~53, 2004.

\bibitem{Neukirch1999}
J.~Neukirch,
\emph{Algebraic Number Theory},
Springer, Berlin, 1999.

\bibitem{Pascadi2025}
A.~Pascadi,
Level of distribution $5/8$ for smooth moduli,
arXiv:2505.00653v2, 2025.

\bibitem{Stark1967}
H.~M.~Stark,
A complete determination of the complex quadratic fields of
class-number one,
\emph{Michigan Math.\ J.} \textbf{14} (1967), 1--27.

\bibitem{Stark1974}
H.~M.~Stark,
Some effective cases of the Brauer--Siegel theorem,
\emph{Invent.\ Math.} \textbf{23} (1974), 135--152.

\bibitem{Watt1995}
N.~Watt,
Kloosterman sums and a mean value for Dirichlet polynomials,
\emph{J.\ Number Theory} \textbf{53} (1995), 179--210.

%% ── Companion papers (Wakil's Program) ───────────────────────────────────────

\bibitem{WakilKhayyam}
K.~Wakil,
Khayyam's triangle: historical restoration and the geometry hidden in the
numbers,
ARC Institute of Knowware, preprint, 2026.

\bibitem{WakilSC}
K.~Wakil,
Spherical cow philosophy: a methodological framework for mathematical
discovery,
ARC Institute of Knowware, preprint, 2026.

\bibitem{WakilShannon}
K.~Wakil,
The Shannon--Wakil effect: constitutional forcing as a universal pattern
in information theory and prime arithmetic,
ARC Institute of Knowware, preprint, 2026.

\bibitem{WakilRamanujan}
K.~Wakil,
Ramanujan's dimensional forcing: a universal formula for sieve levels,
ARC Institute of Knowware, preprint, 2026.

\bibitem{WakilEisenstein}
K.~Wakil,
Ternary forcing on the hexagonal lattice: Eisenstein integers, cascade
moduli, and level of distribution~$5/8$,
ARC Institute of Knowware, preprint, 2026.

\end{thebibliography}

\end{document}
